%% kapa: exam=1 round=21 year=2010 period=2 form=A subject=appraisal_law subject_ko=부동산관계법규 seq=27 number=67
%% source: Q-Net 감정평가사 기출문제 (cst00309) artlSeq=5024181 file=extracted/1008/5024181/2010년 제21회 감정평가사 1차시험 기출문제/2교시.hwp
%% conversion: kordoc 4.24.51 -> tools/gen1.py (auto); review=proofread-gpt-sol; flags=none
\begin{question}{67}
국유재산법령상 국유재산에 관한 설명으로 옳지 않은 것은?
\begin{choices}
  \item 행정재산으로 기부하는 재산에 대하여 기부자가 자신에게 무상으로 사용허가하여 줄 것을 조건으로 기부하는 경우에는 총괄청은 기부를 받을 수 없다.
  \item 총괄청은 관리청이 없거나 분명하지 아니한 국유재산에 대하여 그 관리청을 지정한다.
  \item 일반재산은 대통령령으로 정하는 바에 따라 부동산신탁을 취급하는 신탁업자에게 신탁하여 개발할 수 있다.
  \item 행정재산의 사용허가를 받은 자가 그 행정재산의 관리를 소홀히 하여 재산상의 손해를 발생하게 한 경우에는 사용료 외에 대통령령으로 정하는 바에 따라 그 사용료를 넘지 아니하는 범위에서 가산금을 징수할 수 있다.
  \item 총괄청은 국유재산의 효율적인 관리와 처분을 위하여 필요하다고 인정하면 일정한 범위를 정하여 일반재산을 지방자치단체에 일괄하여 양여할 수 있다.
\end{choices}
\end{question}
\answer{1}
